\chapter{Instructions as Functions}
\label{ch:functions}

An instruction is a function from machine states to machine states. A program is a composition of such functions. Two programs are \term{equivalent} when they compute the same function, and a compiler optimization is a claim that two programs are equivalent and one of them is cheaper.

This chapter is about the first half of that claim, and about why it is the easy half.

\section{Equivalence is a for-all}
\label{sec:equiv}

Take Warren's population count from \emph{Hacker's Delight}, Figure 5-2: five lines of shifts, masks and adds that count the set bits of a thirty-two-bit word. Its specification is the obvious one --- add up the thirty-two bits. The claim that the trick computes the specification is a sentence with one universal quantifier over \(2^{32}\) words, and it is the kind of sentence a solver decides in well under a second.

\begin{artifact}{I.thm.popcount32}
The five-step SWAR population count equals the thirty-two-term bit sum for every thirty-two-bit word. This is the worked example in the Lean 4 reference manual's chapter on bitvectors, proved there by verified bit-blasting; it is the first identity in the 2025 result that decided \emph{all} of Hacker's Delight for every width at once; and it is decided here, at width thirty-two, by the same front door as the Prologue.
\ArtifactScope{I.thm.popcount32}
\end{artifact}

\begin{codelisting}[language=Lisp,caption={The negated equivalence, as the solver sees it. The specification is the thirty-two-term sum; the theorem is that the difference has no model.}]
(set-logic QF_BV)
(declare-const x (_ BitVec 32))
(define-fun step1 () (_ BitVec 32)
  (bvsub x (bvand (bvlshr x #x00000001) #x55555555)))
; ... four more steps ...
(define-fun hd () (_ BitVec 32) (bvand step5 #x0000003F))
(assert (not (= hd spec)))
(check-sat)          ; unsat
\end{codelisting}

Twenty years of verified compilation certify sentences of this shape. CompCert's peephole optimizer already runs the architecture this book advocates --- an untrusted oracle proposes a rewrite, symbolic execution validates it --- and translation validators like Alive2 do the same for LLVM one rewrite at a time. The property they certify is always the same one: the new code computes what the old code computed.

\section{Minimality is a non-existence}
\label{sec:minimality}

Now take the other half. To say that a program is \emph{minimal} is to say that no shorter program in some declared instruction set computes the same function. That is not a for-all over inputs; it is a for-all over \emph{programs}, and its negation --- \emph{there exists a shorter program} --- is a search. Refuting a search is exactly what a resolution proof is for.

\begin{keyidea}[title={The asymmetry in one sentence}]
Correctness is a universal statement discharged once per rewrite. Minimality is a non-existence statement over a whole search space. The first has been certified thousands of times; the second, for an instruction sequence on any real machine, had not been certified by anyone before the objects in Part IV of this book.
\end{keyidea}

\begin{artifact}{I.prin.equivalence-vs-minimality}
Stated here as a principle, with no artifact, because it is not a claim about a machine; it is the claim the rest of the book exists to make good on.
\ArtifactScope{I.prin.equivalence-vs-minimality}
\end{artifact}

It is worth being precise about how absent the second half is, because the field's vocabulary hides it. The superoptimization literature runs from Massalin's 1987 exhaustive search of a 68020 subset --- verified, in his words, by having the author \emph{look at} the output --- through Bansal and Aiken's three-instruction x86 peepholes in 2006, to HieraSynth in 2025, which synthesizes seven- and eight-instruction programs over seven-hundred-instruction sets and states outright that its completeness \emph{proves the generated program is optimal}. Every one of these establishes minimality the same way: the search space is exhausted, a solver says \emph{unsat} at each leaf, and the refutation is discarded. HieraSynth's escape hatch for readers who want more is to \emph{disable timeouts for leaf nodes}. The proof is the solver returning \emph{unsat} with the clock switched off.

Meanwhile the SAT community has required machine-checkable proofs of every \emph{unsat} in its main competition track since 2013, and disqualifies solvers whose certificates fail to check. Both halves of the machinery exist. Nobody had joined them for instruction sequences.

\section{What an instruction set is, in this book}
\label{sec:isa}

Because minimality quantifies over programs, it quantifies over an \emph{instruction set}, and the set has to be a term, not a name. \enquote{AVX2} is a name; it covers several hundred instructions across data movement, arithmetic, memory, masking and zeroing, with operand forms that change the semantics. A minimality claim \enquote{over AVX2} is a claim nobody can check because nobody can say what it quantifies over.

Every theorem in Part IV therefore names its instruction set explicitly --- five families and their operand forms; fourteen families and an SSA operand model --- and the ledger prints that scope in every artifact box. The reader will notice that this makes the theorems narrower than the headline a reviewer might want. That is correct, and it is the point: a narrow theorem with a checkable certificate is a theorem; a broad claim over an undeclared set is a sentence.

\section{Exercises}

\begin{enumerate}
\item Decide, at width thirty-two, that Warren's bit-reversal routine (Figure 7-1) computes the reverse of the bit sum's argument. Add it to the ledger.
\item The 2025 width-independent result decided one hundred percent of Hacker's Delight and twenty-seven percent of the peephole rewrites in LLVM's test suite. Pick one of the remaining seventy-three percent and decide it at a fixed width. What would you need to state it for all widths?
\item State precisely what \enquote{the instruction set of AVX2} would have to mean for the sentence \enquote{byte reversal needs at least two AVX2 instructions} to be a theorem. How many operand forms does \code{vpshufb} alone have?
\end{enumerate}
